\documentclass{article}
\usepackage{amsmath,amsthm}
\usepackage{subcaption}
\usepackage{booktabs}
\usepackage{maleficium-interactive}
\maleficiumtheme{house}
\maleficiumreaderstyle{reader.css}
\newtheorem{theorem}{Theorem}
\newtheorem{lemma}[theorem]{Lemma}
\theoremstyle{definition}
\newtheorem{definition}[theorem]{Definition}
\title{The Playground Paper}
\author{Maleficium}
\date{}
\begin{document}
\maketitle
\begin{abstract}
One source that compiles to a PDF today and is meant to become an
interactive HTML paper later. Every widget macro the package offers
appears at least once, beside live equations, plain figures, subfigures, math,
citations that preview their reference on hover, cross-references, and
theorems whose proofs fold away in the HTML paper.
\end{abstract}

\section{Introduction}
Text typeset with \TeX{}~\cite{knuth1984texbook} and \LaTeX{}~\cite{lamport1994latex}
(together,~\cite{knuth1984texbook,lamport1994latex}) reads the same in print and on screen; the widgets below add an
interactive layer on top of that~\cite{shannon1948mathematical}.
Section~\ref{sec:guide} explains each widget. The model poses in
Figures~\ref{fig:model} and~\ref{fig:model-pose} and the clips play in
Figures~\ref{fig:video} and~\ref{fig:video-remote}; Table~\ref{tab:results}
reads a CSV, Figures~\ref{fig:chart}--\ref{fig:chart-styled} draw charts,
Figures~\ref{fig:plotly-tradeoff} and~\ref{fig:plotly-latency} draw Plotly
figures, and Figures~\ref{fig:html}--\ref{fig:wave} embed HTML pages and a
custom runtime. Plain figures, subfigures, math and theorems follow.

\section{Widgets}\label{sec:widgets}

\subsection{Model}
\begin{figure}[ht]
\centering
\caption{A three-dimensional model widget; the PDF shows a render of the model.}
\label{fig:model}
\interactiveruntime[runtime=model@1, height=5cm, id=fig-model,
  alt={A colored cube, the interactive 3D mesh}]{models/mesh.glb}
\end{figure}

\subsection{Model from a posed camera}
Figure~\ref{fig:model-pose} shows the same cube through a posed camera: a
\texttt{size=} in pixels shapes the poster, a white \texttt{background=}
sits behind the cube, and a \texttt{camera=} puts the eye at $(0,1,3)$
looking at the origin. The PDF shows the rendered poster; the live view
starts from that camera.
\begin{figure}[ht]
\centering
\caption{The cube from the front at eye level, on a white backdrop.}
\label{fig:model-pose}
\interactiveruntime[runtime=model@1, height=5cm, id=fig-model-pose, size=640x480,
  background=ffffff, camera={pos=0,1,3 target=0,0,0 up=0,1,0},
  alt={The colored cube from the front at eye level on a white backdrop}]{models/mesh.glb}
\end{figure}

\subsection{Video}
The clip loops quietly: \texttt{loop=true} plays it again from the start
and \texttt{muted=true} starts it silent. Figure~\ref{fig:video-remote} also
names a hosted copy: \texttt{remote=} with its \texttt{sha256=} is the copy
used when the file does not travel with the paper, and \texttt{at=} starts
the live view (and draws its poster) one second in.
\begin{figure}[ht]
\centering
\caption{A video widget; the PDF shows a frame of the clip.}
\label{fig:video}
\interactiveruntime[runtime=video@1, poster=figures/clip.png, height=5cm, id=fig-video,
  loop=true, muted=true,
  alt={A frame of the test-pattern clip}]{media/clip.mp4}
\end{figure}
\begin{figure}[ht]
\centering
\caption{The test-pattern clip from its remote copy, starting one second in.}
\label{fig:video-remote}
\interactiveruntime[runtime=video@1, height=4cm, id=fig-video-remote, at=1,
  remote=https://example.org/media/clip.mp4,
  sha256=20ae41ca929a117aa92fd403bff6fd7734b5d97ff223d91f7873d7d62050f738,
  alt={The test-pattern clip, starting one second in}]{media/clip.mp4}
\end{figure}

\subsection{Table from CSV}
Table~\ref{tab:results} is read from a CSV file. The PDF keeps the first
three rows; the interactive view starts sorted by accuracy, filtered to the
\texttt{ours} rows, showing three of them with two decimals.
\begin{table}[ht]
\centering
\caption{Accuracy and latency for each configuration.}
\label{tab:results}
\interactiveruntime[runtime=table@1, pdfrows=3, id=tab-results,
  sort={accuracy desc}, filter=ours, columns={method;accuracy;latency},
  format={accuracy=fixed:2;latency=fixed:0}, maxrows=3,
  alt={Accuracy and latency for five configurations, from baseline to ours-xl}]{data/results.csv}
\end{table}

\subsection{Chart from a Vega-Lite spec}
Figure~\ref{fig:chart} draws accuracy by method from a Vega-Lite
specification; \texttt{scale=2} renders its poster at twice the pixels per
CSS pixel.
\begin{figure}[ht]
\centering
\caption{Accuracy by method, drawn from a Vega-Lite specification.}
\label{fig:chart}
\interactiveruntime[runtime=chart@1, height=5cm, id=fig-chart, scale=2,
  alt={Bar chart of accuracy rising from 0.71 for baseline to 0.89 for ours-xl}]{charts/ablation.vl.json}
\end{figure}

\subsection{Chart from two data files}
Figure~\ref{fig:chart-split} layers accuracy and latency from two CSV files:
a chart's \texttt{data} source may repeat, and the spec names each file it
reads.
\begin{figure}[ht]
\centering
\caption{Accuracy and latency by method, each from its own CSV file.}
\label{fig:chart-split}
\interactiveruntime[runtime=chart@1, height=5cm, id=fig-chart-split,
  sources={data=data/accuracy.csv, data=data/latency.csv},
  alt={Layered bars of accuracy and latency by method, from two CSV files}]{charts/compare.vl.json}
\end{figure}

\subsection{Chart with its own poster and plate}
Figure~\ref{fig:chart-styled} carries the author's own poster, a narrower
\texttt{width=}, and a \texttt{plate=} of its own instead of the page
colour.
\begin{figure}[ht]
\centering
\caption{Accuracy by method, with its own poster and plate.}
\label{fig:chart-styled}
\interactiveruntime[runtime=chart@1, height=5cm, width=0.6\linewidth, plate=FFF4E0,
  poster=figures/chart-poster.png, id=fig-chart-styled,
  alt={Bar chart of accuracy rising from 0.71 for baseline to 0.89 for ours-xl}]{charts/ablation.vl.json}
\end{figure}

\subsection{Plotly figures from Python}
Figures~\ref{fig:plotly-tradeoff} and~\ref{fig:plotly-latency} are what
Python's \texttt{fig.write\_json()} saves, written from the same CSV by
\texttt{tools/plotly\_figures.py}. The first keeps Plotly Express defaults,
so it takes the paper's style, and its poster renders at \texttt{scale=2}.
The second brings its own template, and
\texttt{template=figure} keeps it.
\begin{figure}[ht]
\centering
\caption{Accuracy against latency, a scatter plot in the paper's style.}
\label{fig:plotly-tradeoff}
\interactiveruntime[runtime=plotly@1, height=5cm, id=fig-plotly-tradeoff, scale=2,
  alt={Line of accuracy against latency, from baseline at 12 ms to ours-xl at 58 ms}]{charts/tradeoff.plotly.json}
\end{figure}
\begin{figure}[ht]
\centering
\caption{Latency by method, a bar chart in the figure's own style.}
\label{fig:plotly-latency}
\interactiveruntime[runtime=plotly@1, template=figure, height=5cm, id=fig-plotly-latency,
  alt={Horizontal bars of latency, from 12 ms for baseline to 58 ms for ours-xl}]{charts/latency.plotly.json}
\end{figure}

\subsection{HTML widgets}
\begin{figure}[ht]
\centering
\caption{A self-contained HTML widget in a sandboxed frame.}
\label{fig:html}
\interactive[height=5cm, id=fig-html,
  alt={The live HTML demo}]{widgets/demo/}
\end{figure}

Figure~\ref{fig:bridge} speaks the bridge: it announces readiness and
status and answers snapshots, drawing the paper's tokens so the reader
stylesheet reaches it. Figure~\ref{fig:framed} declares the origins it may
frame or load from, in its \texttt{widget.json} and in the macro's
\texttt{framedomains=} and \texttt{resourcedomains=} (the two join); it
fetches nothing, and its PDF poster is a placeholder until the widget is
approved and Maleficium renders it.
\begin{figure}[ht]
\centering
\caption{An HTML widget that answers the bridge, drawn from the paper's tokens.}
\label{fig:bridge}
\interactive[poster=figures/bridge-poster.png, height=4cm, id=fig-bridge,
  alt={Swatches of the paper accent and four palette steps}]{widgets/bridge/}
\end{figure}
\begin{figure}[ht]
\centering
\caption{An HTML widget declaring the origins it may frame or load from.}
\label{fig:framed}
\interactive[height=3cm, id=fig-framed, framedomains=https://example.org,
  resourcedomains=https://cdn.example.org,
  alt={A widget declaring frame and resource origins}]{widgets/framed/}
\end{figure}

\subsection{Custom runtime}
Figure~\ref{fig:wave} runs a custom runtime: the waving robot from the
author's website as a real runtime with its own \texttt{runtime.json}. It
speaks the bridge, takes the paper's figure tokens, and reads its speech
bubble from a greeting file; the runtime is not approved in a fresh
checkout, so the PDF holds a placeholder where its proposed poster will go.
\begin{figure}[ht]
\centering
\caption{A waving robot, a custom runtime; its poster arrives with approval.}
\label{fig:wave}
\interactiveruntime[runtime=wave@1, height=5cm, width=0.75\linewidth, plate=1A1233,
  id=fig-wave,
  alt={A pixel-art robot waving and watching the mouse}]{words/greeting.txt}
\end{figure}

\subsection{Widgets side by side}
Figure~\ref{fig:pair} holds one figure with two subfigures, each embedding
its own widget: the left panel reuses the chart spec as
Figure~\ref{fig:pair-a}, the right panel as Figure~\ref{fig:pair-b}.
\begin{figure}[ht]
\centering
\begin{subfigure}{0.45\textwidth}
\centering
\interactiveruntime[runtime=chart@1, height=3cm, id=fig-pair-a,
  alt={Small copy of the accuracy chart, left panel}]{charts/ablation.vl.json}
\caption{Left panel.}\label{fig:pair-a}
\end{subfigure}\hfill
\begin{subfigure}{0.45\textwidth}
\centering
\interactiveruntime[runtime=chart@1, height=3cm, id=fig-pair-b,
  alt={Small copy of the accuracy chart, right panel}]{charts/ablation.vl.json}
\caption{Right panel.}\label{fig:pair-b}
\end{subfigure}
\caption{Two widgets side by side.}
\label{fig:pair}
\end{figure}

\subsection{Widget in a float without a caption}
A float needs no caption. This one holds a widget and nothing else, so it
carries no label and no figure number.
\begin{figure}[ht]
\centering
\interactiveruntime[runtime=chart@1, height=3cm, id=chart-bare,
  alt={Small uncaptioned copy of the accuracy chart}]{charts/ablation.vl.json}
\end{figure}

A widget may also sit directly in the text, outside any float. It records
no label and no figure number either.

\begin{center}
\interactiveruntime[runtime=chart@1, height=3cm, id=chart-inline,
  alt={Small inline copy of the accuracy chart outside any float}]{charts/ablation.vl.json}
\end{center}

\section{Plain floats}
Not every figure is interactive. Figure~\ref{fig:plain} has a caption;
the one after it has none.
\begin{figure}[ht]
\centering
\includegraphics[height=3cm]{figures/plain.png}
\caption{A static figure with a caption.}
\label{fig:plain}
\end{figure}
\begin{figure}[ht]
\centering
\includegraphics[height=3cm]{figures/plain.png}
\end{figure}

\subsection{Subfloats}
Figure~\ref{fig:subs} uses the standard \texttt{subcaption} package, with
Figures~\ref{fig:suba} and~\ref{fig:subb} as its panels.
\begin{figure}[ht]
\centering
\begin{subfigure}{0.45\textwidth}
\centering
\includegraphics[width=\linewidth]{figures/left.png}
\caption{First panel.}\label{fig:suba}
\end{subfigure}\hfill
\begin{subfigure}{0.45\textwidth}
\centering
\includegraphics[width=\linewidth]{figures/right.png}
\caption{Second panel.}\label{fig:subb}
\end{subfigure}
\caption{Two subfigures side by side.}
\label{fig:subs}
\end{figure}

\section{Math and theorems}
Entropy is defined in Equation~\eqref{eq:entropy} and bounded in
Theorem~\ref{thm:bound}. In the HTML paper, hovering or focusing a
citation shows its reference, and each proof can be folded away.
\begin{definition}[Entropy~\cite{shannon1948mathematical}]\label{def:entropy}
The entropy of a random variable $X$ taking $n$ values with probabilities
$p_1, \dots, p_n$ is
\begin{equation}\label{eq:entropy}
H(X) = -\sum_{i=1}^{n} p_i \log p_i .
\end{equation}
\end{definition}
\begin{lemma}[Gibbs' inequality]\label{lem:gibbs}
For distributions $p$ and $q$ on the same $n$ outcomes,
$-\sum_i p_i \log p_i \le -\sum_i p_i \log q_i$~\cite[Thm.~2.6.3]{cover2006elements}.
\end{lemma}
\begin{proof}
Since $\log x \le x - 1$ for every $x > 0$, with equality only at $x = 1$,
\begin{align*}
\sum_i p_i \log \frac{q_i}{p_i}
  &\le \sum_i p_i \left( \frac{q_i}{p_i} - 1 \right) \\
  &= \sum_i q_i - \sum_i p_i \le 0 . \qedhere
\end{align*}
\end{proof}
\begin{theorem}\label{thm:bound}
For a distribution on $n$ outcomes, $0 \le H(X) \le \log n$.
\end{theorem}
\begin{proof}[Proof of Theorem~\ref{thm:bound}]
Each term of Equation~\eqref{eq:entropy} is nonnegative, so $H(X) \ge 0$.
Lemma~\ref{lem:gibbs} with the uniform $q_i = 1/n$ gives
$H(X) \le -\sum_i p_i \log \tfrac{1}{n} = \log n$.
\end{proof}

\subsection{Live equations}
A \verb|liveequation| is an equation whose parameters the reader can
change. The PDF shows the values at their defaults; the HTML paper puts a
slider and a number box under the equation and recomputes each result.
Equation~\eqref{eq:binary-entropy} is Definition~\ref{def:entropy} for two
outcomes with probabilities $p$ and $1-p$, in bits.
\begin{liveequation}[id=eq-binary-entropy,
    alt={Binary entropy in bits at a probability p the reader can change}]
  H_2(\maleficiumparam{p}{0.25}[0.01,0.99,0.01])
    = \liveresult[digits=3]{-(p*ln(p) + (1-p)*ln(1-p))/ln(2)}\ \text{bits}
  \label{eq:binary-entropy}
\end{liveequation}
A result may come before the parameters it uses, a parameter may be
shown again, and \verb|liveequation*| is unnumbered. With
\verb|mode=widget| the equation runs in a frame of its own, on the
figure plate, instead of in the article's text:
\begin{liveequation*}[id=eq-parabola, mode=widget,
    alt={A parabola y = a x squared with live a and x}]
  y = \liveresult{a*x^2}
    = \maleficiumparam{a}{1}[0.5,1.5,0.1] \cdot \maleficiumparam{x}{1.5}[-2,2,0.1]^2,
  \qquad x = \maleficiumparam{x}
\end{liveequation*}

\section{What each widget does}\label{sec:guide}
\begin{description}
\item[Model] \verb|runtime=model@1| loads a glTF model (\texttt{.glb}) for orbiting
  and zooming. The PDF shows the poster (Figure~\ref{fig:model}). A posed
  camera starts the live view from \verb|camera=|, with a \verb|size=| poster
  in pixels and a \verb|background=| behind the model (Figure~\ref{fig:model-pose}).
\item[Video] \verb|runtime=video@1| plays a clip, with an optional remote URL
  and checksum for large files; \verb|loop=true| plays it again and
  \verb|muted=true| starts it silent (Figure~\ref{fig:video}). With
  \verb|at=| the live view starts seconds in and the poster is that frame,
  and \verb|remote=| with \verb|sha256=| names the hosted copy used when the
  file does not travel with the paper (Figure~\ref{fig:video-remote}).
\item[Table] \verb|runtime=table@1| reads a CSV into a sortable table. The PDF
  prints only the first \texttt{pdfrows} rows (Table~\ref{tab:results}); the
  live view starts sorted by \verb|sort=|, filtered by \verb|filter=|, showing
  the \verb|columns=| in order with the \verb|format=| each names and at most
  \verb|maxrows| rows.
\item[Chart] \verb|runtime=chart@1| renders a Vega-Lite specification with
  hover details, its poster at \verb|scale=| pixels per CSS pixel
  (Figure~\ref{fig:chart}). A chart's \verb|data| source may repeat, with the
  spec naming each file it reads (Figure~\ref{fig:chart-split}); an author
  poster, a narrower \verb|width=| and a \verb|plate=| of its own style one
  widget (Figure~\ref{fig:chart-styled}).
\item[Plotly] \verb|runtime=plotly@1| draws a figure saved from Python with
  Plotly's hover, in the paper's style unless \verb|template=figure|
  (Figures~\ref{fig:plotly-tradeoff} and~\ref{fig:plotly-latency}); its
  poster also takes \verb|scale=|.
\item[HTML] \verb|\interactive| embeds a folder with an \texttt{index.html} in a
  sandboxed frame (Figure~\ref{fig:html}). One page answers the bridge, with
  readiness, status and snapshots drawn from the paper's tokens
  (Figure~\ref{fig:bridge}); another declares the origins it may frame or
  load from, in its \texttt{widget.json} with the macro's
  \verb|framedomains=| and \verb|resourcedomains=| (Figure~\ref{fig:framed}).
\item[Custom runtime] \verb|\interactiveruntime| with a project runtime runs
  the project's own page under the bridge: readiness, status, snapshots and
  the paper's tokens, with its sources bound from files
  (Figure~\ref{fig:wave}).
\item[Live equation] \verb|liveequation| with \verb|\maleficiumparam{name}{default}[min,max,step]|
  and \verb|\liveresult[digits=n]{formula}|: the PDF typesets the values at the
  defaults, the HTML paper lets the reader change them
  (Equation~\eqref{eq:binary-entropy}). A widget may also sit in the text
  outside any float, at the author's width, with no number; the live equation
  runs there by default and in a frame of its own with \verb|mode=widget|.
\item[Reader style] \verb|\maleficiumreaderstyle{reader.css}| restyles the
  reader in light and dark mode alike; its tokens reach every widget and the
  posters Maleficium renders (the bridge widget's accent chip shows them).
\end{description}
Every widget requires \texttt{alt=} text, which becomes the image
description in the HTML view. The app renders the PDF's picture of a
model, chart, Plotly figure or HTML widget from the widget itself, so these carry no
\texttt{poster=}; a video takes a \texttt{poster=} frame or an \texttt{at=}
moment, and a widget whose poster needs an approval the checkout does not
have yet shows a placeholder (Figures~\ref{fig:framed} and~\ref{fig:wave}).

\bibliographystyle{plain}
\bibliography{refs}
\end{document}
